% Avsnitt 19.9, ur lectures/L19/README.md.

\section{Sammanfattning}
\label{c19:sec:review}

\needspace{6\baselineskip}
\subsection*{Det här ska du kunna nu}
\begin{itemize}
  \item Implementera transaktions-FSM:en enligt protokollspecifikationen, inklusive avbrottsregeln.
  \item Förklara varför en läsning låser registervärdet en gång, och vad som går fel utan den regeln.
  \item Komponera hela noden till en toppnivå och verifiera den i simulering.
  \item Koppla ihop två DE0-CV till en delad buss, och en mikrokontroller som master över dess
    MVIO-domän, och felsöka kedjan systematiskt nerifrån och upp.
  \item Säga vilka krav SPI-kopplingen ställer på hattens konstruktion, och varför de inte går att
    rätta i efterhand.
  \item Säga vad simulering har bevisat och vad bara hårdvara kan bevisa.
\end{itemize}

\needspace{6\baselineskip}
\subsection*{Frågor att testa dig själv med}
\begin{enumerate}
  \item Varför låses registervärdet en gång vid slutet av kommandobyten i stället för att läsas medan
    databytesen skiftas ut? Vilket register lider mest av att den regeln bryts?
  \item SS går hög efter tre bytes. Vad har hänt med transaktionen, och vad har ändrats i noden?
  \item En skrivning till ett reserverat index. Vad gör bryggan, och varför konsumerar den ändå alla
    fem bytesen?
  \item Varför finns ingen felkanal i transaktionslagret?
  \item Vilket bring-up-steg skulle avslöja en förväxlad MISO och MOSI, och vilket skulle avslöja ett
    \code{VDDIO2} som aldrig kopplades in?
  \item Hatten lägger en lysdiod på \code{PC2}. Vilket steg fallerar först, och varför är det ett fel
    som inte går att programmera sig ur?
  \item Vad bevisade stegen 1 till 3 som stegen 4 till 8 inte behövde bevisa om igen?
  \item Vad har simuleringen bevisat om designen, och vad har den inte kunnat bevisa?
\end{enumerate}

\needspace{6\baselineskip}
\subsection*{Slutredovisning}
Varje grupp visar, på ungefär femton minuter: en körning av \code{make build-project} med alla
testbänkar gröna; en kort genomgång av \code{can_spi_node}, vad varje block gör och var gränserna
går; bring-upen, så långt ni kommit, på riktig hårdvara; och \file{CONTRIBUTORS.md} plus en snabb
genomgång av repots historik, branchar, PR:er och granskningar.

Slutredovisningen är projektets andra och sista fasta hållpunkt. Projektet lämnas in samma dag; se
checklistan i bilaga~\ref{app:project}.

Bring-upen bedöms inte. Kommer du till steg 3 av åtta är det ett bra resultat; det som bedöms är
simuleringen bakom den.

\needspace{6\baselineskip}
\subsection*{Blicken framåt}
\Chapref{c20:ch} är bokens andra praktiska tentamen, med uppgifter på projektets moduler. Efter den är
kursen slut.

Kontrollern du byggt lämnas över till parallellklassen, som skriver C++-drivrutinen mot samma
register över samma SPI-transport. Det gemensamma kontraktet är bilaga~\ref{app:protocol}, och
ingenting annat.
